\section{One Vector in a Different Basis}

A vector does not change when we describe it in another basis. Only the numbers
used to reconstruct it from the chosen basis vectors change. One numerical
example is enough to see the mechanism.

Let
\[
\mathbf b_1=[1,1],
\qquad
\mathbf b_2=[1,-1]
\]
and consider the vector
\[
\mathbf v=[4,2].
\]
We seek numbers $\alpha$ and $\beta$ such that
\[
\mathbf v=\alpha\mathbf b_1+\beta\mathbf b_2.
\]
Using the Cartesian coordinates of the basis vectors gives
\[
[4,2]
=
\alpha[1,1]+\beta[1,-1]
=
[\alpha+\beta,\alpha-\beta].
\]
Therefore
\[
\alpha+\beta=4,
\qquad
\alpha-\beta=2.
\]
Solving this small system gives
\[
\alpha=3,
\qquad
\beta=1.
\]
Hence
\[
\boxed{
\mathbf v=3\mathbf b_1+\mathbf b_2
}
\]
and the coordinates of the same vector in the basis
$B=(\mathbf b_1,\mathbf b_2)$ are
\[
\boxed{
[\mathbf v]_B=[3,1].
}
\]

\begin{center}
\begin{tikzpicture}[scale=0.9]
    \coordinate (O) at (0,0);
    \coordinate (BOne) at (1.4,1.4);
    \coordinate (BTwo) at (1.4,-1.4);
    \coordinate (ThreeBOne) at (4.2,4.2);
    \coordinate (V) at (5.6,2.8);

    \draw[subset,->] (O) -- (BOne)
        node[midway,above left] {$\mathbf b_1$};
    \draw[subset,->] (O) -- (BTwo)
        node[midway,below left] {$\mathbf b_2$};
    \draw[helper,dashed,->] (O) -- (ThreeBOne)
        node[midway,above left] {$3\mathbf b_1$};
    \draw[helper,dashed,->] (ThreeBOne) -- (V)
        node[midway,above right] {$\mathbf b_2$};
    \draw[very thick,->] (O) -- (V)
        node[midway,below right] {$\mathbf v$};

    \fill (O) circle (2pt) node[below left] {$O$};
    \fill (ThreeBOne) circle (2pt);
    \fill (V) circle (2pt);
\end{tikzpicture}
\end{center}

The Cartesian coordinates $[4,2]$ and the basis coordinates $[3,1]_B$ are two
numerical descriptions of the same geometric vector. Changing the basis means
solving for the coefficients in one concrete linear combination; no additional
general machinery is needed here.
